\chapter{Common emitter amplifier}
This experiment explores the usage of a transistor as an amplifier.

\section{Theory}\label{sec:theory}
\begin{figure}[tbp]
	\centering
	%\includegraphics[width=.5\textwidth]{./img/com_emitter_amp.pdf}
	\caption[Single Stage Common Emitter Amplifier]{Single stage common emitter amplifier}
	\label{fig:com_emitter_amp}
\end{figure}
The circuit shown in \autoref{fig:com_emitter_amp} acts as an inverting voltage amplifier.
The input and output capacitors remove any DC component of the input.
\comp{R_1} and \comp{R_2} are used to bias the transistor, so that it remains in active mode during operation.
This method of biasing vastly reduces the effects of varying $\beta$ as a cause of temperature drift by holding the base bias at a constant voltage.
Adversely, the constant voltage at the base introduces a stronger dependency on the base-emitter voltage drop, which changes by \SI{-2}{\milli\volt\per\kelvin}.
\comp{R_E} is used to compensate for this by introducing negative feedback, albeit lowering the gain.

Temperature drift can be calculated as
\begin{equation}
	\frac{\d U_\text{a}}{\d T}=\frac{\d U_\text{a}}{\d U_\text{e}}\frac{\d U_\text{e}}{\d T}\approx A\cdot\SI{-2}{\milli\volt\per\kelvin},
\end{equation}

where $U_\text{a}$ denotes the output voltage, $U_\text{e}$ the input voltage and $T$ the temperature.
This means that temperature drift is proportional to the gain $A$ in a common emitter amplifier.
As temperature drift happens slowly, only DC amplification has to be decreased, whereas AC amplification can remain unchanged.
This can be realized by putting a capacitor \comp{C_E} in parallel to the emitter resistor \comp{R_E} which will decrease the negative feedback with growing frequency.

The quiescent base voltage is determined by the voltage divider formed by \comp{R_1} and \comp{R_2}:
\begin{equation}\label{eq:bias_volt}
	U_\text{bias}=U_\text{cc}\cdot\frac{\comp{R_2}}{\comp{R_1}+\comp{R_2}}.
\end{equation}

\begin{figure}[tbp]
	\centering
	%\includegraphics[width=.5\textwidth]{./img/small_signal_neg_feed.pdf}
	\caption{Small signal schematic of a common emitter circuit with negative feedback}
	\label{fig:small_signal_neg_feed}
\end{figure}

To deduce a proper formula for the AC-voltage amplification of a common emitter circuit with negative feedback, let us consider the small signal equivalent schematic seen in \autoref{fig:small_signal_neg_feed}.
The transistor can be replaced by a constant current source with differential resistances \comp{r_{BE}} and \comp{r_{CE}}.
Using \texttt{Kirchhoff}'s law for nodes, it holds
\begin{alignat}{2}
	\frac{u_\text{e}-u_\text{E}}{r_\text{BE}}+Su_\text{BE}+\frac{u_\text{a}-u_\text{E}}{r_\text{CE}} &= \frac{u_\text{E}}{R_\text{E}} \label{eq:small_1}\\
	Su_\text{BE}+\frac{u_\text{a}-u_\text{E}}{r_\text{CE}} + \frac{u_\text{a}}{R_\text{C}}&= i_\text{a}, \label{eq:small_2}
\end{alignat}
where $S=\frac{I_\text{C}}{U_\text{T}}$ denotes the transistor's transconductance.
Rearraging equations \ref{eq:small_1}, \ref{eq:small_2} and using $u_\text{BE}=u_\text{e}-u_\text{E}$, this yields
\begin{align}
	A&=\frac{u_\text{a}}{u_\text{e}}\at[\bigg]{i_\text{a}=1} \nonumber \\
	&=-\frac{SR_\text{C}\left(1-\frac{R_\text{E}}{\beta r_\text{CE}}\right)}{1+R_\text{E}\left(S\cdot\left(1+\beta^{-1}+\frac{R_\text{C}}{\beta r_\text{CE}}\right)+r_\text{CE}^{-1}\right)+\frac{R_\text{C}}{r_\text{CE}}} \nonumber \\
	&\approx -\frac{SR_\text{C}}{1+SR_\text{E}}, \label{eq:amp_neg_fb}
\end{align}
where $r_\text{CE}\gg R_\text{C},R_\text{E}$, $\beta\gg 1$ and $\beta$ denotes its current amplification $\frac{I_\text{c}}{I_\text{B}}$.
\section{Evaluation}
\begin{figure}[tbp]
	\centering
	%\includegraphics[width=.5\textwidth]{./img/com_emitter_in_out.pdf}
	\caption{Input versus output voltage}
	\label{fig:com_emitter_in_out}
\end{figure}
\begin{table}[b!]
	\centering
	\caption{Input, output voltages $U_\text{pp,in/out}$ and resulting voltage amplification $A$ at $f=\SI{1}{\kilo\hertz}$}
	\label{tab:com_emit_in_out}
	\begin{tabular}{S|SSS}
		\toprule
		{}&	{$U_\text{pp,in}$ in $\si{\milli\volt}$}&	{$U_\text{pp,out}$ in $\si{\volt}$}&	{$A$}\\
		\midrule
		{with \comp{C_e}}	&	\num{19}	&	\num{2.7}	&	\num{-142}\\
		{}	&	\num{68}	&	\num{10.3}	&	\num{-151}\\
		\midrule
		{without \comp{C_e}}	&	\num{70.8}	&	\num{0.322}	&	\num{-4.55}\\
		{}	&	\num{278}	&	\num{1.3}	&	\num{-4.676}\\
		\bottomrule
	\end{tabular}
\end{table}
The circuit \autoref{fig:com_emitter_amp} is built (including the emitter capacitor, which is taken out later on) and a triangular input signal is applied.
\autoref{fig:com_emitter_in_out} shows the measured voltage curves for input and output voltage.
Minimal non-linearities in the amplified output signal can be observed.
\autoref{tab:com_emit_in_out} lists measured peak-to-peak voltages and resulting voltage amplifications.
Using \autoref{eq:bias_volt} the bias voltage can be calculated as $U_\text{bias}=\SI{2.27}{\volt}$.
Assuming that the base current is negligable, the collector current is $I_\text{c}\approx\frac{U_\text{bias}-\SI{0.7}{\volt}}{\comp{R_E}}=\SI{15.7}{\milli\ampere}$.
Using $I_\text{c}$, we can calculate the conductance to be $S=\frac{I_\text{c}}{U_\text{T}}=\SI{604}{\milli\siemens}$.
The transconductance can then be utilized to finally calculate the theoretical voltage amplification \textbf{without capacitor $\comp{C_E}$} by using the formula derived in \autoref{eq:amp_neg_fb}.
\begin{align*}
	A &= -\frac{SR_\text{C}}{1+R_E} \\
	&=\num{-4.63}.
\end{align*}
With \comp{C_E}, however, the amplification is expected to rise further up to $A=-S\cdot R_\text{C}=\num{283.88}$, which corresponds to the AC-amplification value of a common emitter circuit without negative feedback.
The cutoff frequency of the high-pass formed by \comp{R_E} and \comp{C_E} can be calculated by
\begin{align*}
	f_\text{cutoff}&=\frac{1}{2\pi R_\text{E}C_\text{E}}	\\
	&=\SI{6.37}{\hertz}.
\end{align*}
This means that for frequencies vastly over $f_\text{cutoff}$ negative feedback does not apply anymore.
\begin{figure}[tbp]
	\centering
	\includegraphics[width=.7\textwidth]{./data/plots/1_4.pdf}
	\caption{Amplification $A$ over Frequency $f$}
	\label{fig:com_emitter_freq}
\end{figure}
\begin{table}[b!]
	\centering
	\caption{Input, output voltages $U_\text{in/out}$ at different frequencies $f$}
	\label{tab:com_emit_freq}
	\begin{tabular}{S|SS|SS}
		\toprule
		{Frequency in $\si{\kilo\hertz}$}&
		{$U_\text{in}$ in $\si{\milli\volt}_\text{pp}$}&
		{$U_\text{out}$ in $\si{\volt}_\text{pp}$}&
		{$U_\text{in,cap}$ in $\si{\milli\volt}_\text{pp}$}&
		{$U_\text{out,cap}$ in $\si{\volt}_\text{pp}$}\\
		\midrule
		0.01    &  368  &   0.316  & 40   &   0.4	\\
		0.025   &   416   &  0.688 &  50   &   1	\\
		0.05   &   416   &  1.2   &  51   &  1.66	\\
		0.1  &   416   &  1.7   &  51   &   2.6	\\
		0.5  &   416   &  2     &  50   &   7	\\
		1  &   416   &  2     &  47   &   8	\\
		5  &   416   &  2     &  47   &   8.5	\\
		10 &   408  & 1.96  &  48   &   8.6	\\
		50 &   432  &   2.10  &  47   &   8.3	\\
		100  & 432   &  2.12  &  47   &   8	\\
		\bottomrule
	\end{tabular}
\end{table}

To examine this statement, a sinusoidal signal of varying frequency is applied and input and output voltages are measured to determine the voltage amplification $A$.
\autoref{fig:com_emitter_freq} shows the amplifcation calculated from \autoref{tab:com_emit_freq} over frequency $f$.
As expected, without \comp{C_E} the amplification remains almost constant.

With \comp{C_E}, however, the amplification gradually increases until it reaches its maximum at $\approx\num{176}$, which deviates from the theoretical value of \num{283.88}.
This result may be explained by the ESR (\textbf{E}quivalent \textbf{S}eries \textbf{R}esistance) of \comp{C_E} which poses a residual resistance at high frequencies instead of the capacitor being a short circuit.
Furthermore, the conductance of the transistor in a common emitter circuit without negative feedback differs from the conductance in the examined circuit in that the base current changes with \comp{R_E} not being present.

The highpass frequency response function
\begin{equation*}
	A=\frac{A_0}{\sqrt{1 + (2\pi f\tau)^{-2}}}
\end{equation*}
is fitted to the data with \comp{C_E}, as explained in \autoref{sec:theory}.
A fit parameter of $\tau=\SI{481.3e-6}{\second}^{-1}$ is determined, which deviates from the theoretical value of $RC=\SI{25e-3}{\second}^{-1}$ by \SI{98.1}{\percent}.
The exact frequency response is influenced by many factors, such as time constants formed by the AC-coupling capacitors.
See \url{http://whites.sdsmt.edu/classes/ee320/notes/320Lecture23.pdf} (p.13).
